\chapter*{ABSTRACT}
\addcontentsline{toc}{chapter}{ABSTRACT}
Interior design is difficult for non-professional users because a useful design result must connect the user's design intent with actual room geometry, furniture dimensions, openings, circulation, and visual appearance. Existing solutions address this need only partially. Image generation tools can quickly create attractive interior pictures, but the results are usually static and may not match the actual room layout. Research systems for room arrangement can suggest furniture placements, but their outputs still need geometric validation and editability before they can be used as practical design results.

This thesis proposes RoomFlow, an AI-assisted interior design system for living rooms and bedrooms. The main idea is to keep the editable layout as the central representation of the design, instead of generating a final image directly from text. RoomFlow first generates an initial furniture layout from a user-provided design specification, including the room type, room boundary, fixed openings, furniture inventory, and natural-language design brief. A language-model-based planning pipeline interprets the user's intent and extracts furniture requirements and spatial preferences. A geometric solver then checks containment, collisions, opening clearance, and circulation before returning editable layout variants.

After a layout is selected, RoomFlow reconstructs it in a 2D/3D workspace and uses the chosen 3D view as guidance for layout-guided photorealistic rendering. The rendering process uses the 3D snapshot, visible-object information, and preservation constraints to create a realistic perspective image while preserving the room shell, openings, camera framing, and major furniture arrangement. The system also supports localized image editing, allowing users to add, remove, replace, or recolor selected elements while keeping unrelated parts of the room stable.

The thesis makes three main contributions: a layout design workflow that combines language-based planning with deterministic geometric validation, a layout-guided photorealistic rendering and localized editing workflow, and an integrated web application that connects room creation, layout generation, 3D preview, rendering, and editing. The evaluation shows that RoomFlow improves layout usability, supports editable design exploration, and helps generated images remain consistent with the selected layout.

\begin{flushright}
\begin{minipage}{5cm}
\centering
\textbf{Student}\\[0.1cm]
\includegraphics[width=4cm]{figures/signatures/student_signature_clean.png}\\[-0.35cm]
\studentname
\end{minipage}
\end{flushright}
\clearpage
